% commentaire.tex — f171v, vv. 765–799 (conversion initiale des travaux antérieurs)
\comm{767}{ἐσκυλίσας ἐς τοιαυτὶ πράγματα.}{Pour la métaphore de «~rouler dans les malheurs~», voir le vers 650 ainsi que ce qu'en dit Taillardat (§558).}
\comm{770}{ἐκ τοῦ Παλαμήδους·}{Euripide représenta son \italique{Palamède} en 415 av. J.-C. dans une tétralogie comportant aussi l'\italique{Alexandros}, les \italique{Troyennes} et le drame satyrique \italique{Sysiphe} (il fut classé second). Eschyle et Sophocle représentèrent aussi un \italique{Palamède}. Ce héros aux-mille-tours ingénieux de l'épopée troyenne (et pourtant absent de l'\italique{Iliade} ; sans doute pour ne pas faire d'ombre à Ulysse) avait procuré à l'armée grecque d'« immenses bienfaits »\nc{On retrouve dans une scholie à son \italique{Oreste}, au vers 432, les grandes lignes de la version euripidéenne du drame. J'en reprends la traduction de Jouan \& van Looy, à retrouver dans leur édition des fragments d'Euripide aux Belles Lettres, tome VIII, \nieme{2} partie.} : il apaisa une famine à Aulis et des troubles causés à propos du partage des provisions, en enseignant aux Grecs les « lettres phéniciennes », leur apprit le jeu de dés pour se distraire, les mesures et les jetons de vote, ce qui lui valut un grand renom. Le scholiaste de l'\italique{Oreste} (v. 432) continue :\citComm{ἐπὶ τούτῳ δὲ φθονήσαντες οἱ περὶ Ἀγαμέμνονα καὶ Ὀδυσσέα καὶ Διομήδην τοιόνδε τι σκευωροῦσι κατ᾽ αὐτοῦ. λαβόντες γὰρ Φρύγα αἰχμάλωτον χρυσίον κομίζοντα Σαρπηδόνι ἠνάγκασαν γράψαι Φρυγίοις γράμμασι περὶ προδοσίας ὡς παρὰ Πριάμου πρὸς Παλαμήδην. καὶ τοῦτον μὲν φονεύουσι, θεράποντα δὲ Παλαμήδους πείθουσι χρήμασιν ἅμα τοῖς Τρωϊκοῖς χρήμασι καὶ τὸ γραφὲν πινάκιον ὑπὸ τὴν κλίνην θέσθαι Παλαμήδους. αὐτοὶ δὲ παρελθόντες προδοσίαν κατήγγελλον τοῦ ἥρωος καὶ φωραθῆναι τὴν σκηνὴν ἐκέλευον. εὑρεθέντος δὲ τοῦ πινακίου καὶ τῶν χρημάτων ὑπὸ τὴν κλίνην λίθοις φονεύεται Παλαμήδης\\
\francais{Devenus jaloux de ce fait, Agamemnon, Ulysse et Diomède imaginèrent ce genre de procédé : ayant fait prisonnier un Phrygien qui apportait de l'or à Sarpédon, ils le forcèrent à écrire en caractères phrygiens une lettre parlant de trahison, comme si elle avait été adressée par Priam à Palamède, puis ils le tuèrent. Ils corrompirent alors un serviteur de Palamède, pour qu'il place en même temps que l'or phrygien la tablette écrite sous la tente de son maître. Eux-mêmes, ayant comparu [dans l'assemblée des chefs], accusèrent le héros de trahison et demandèrent qu'on fouille la tente. Après la découverte de la tablette et de l'or sous la tente, Palamède fut lapidé.}}Anxieux de faire connaître la mort injuste de son fils à Nauplios, Œax, le frère de Palamède, imaginait le stratagème ici repris par Mnésiloque\nc{Scholie au vers 771 des \italique{Thesmophories}.} :
\citComm{ὁ γὰρ Ἐὐριπίδης ἐν τῷ Παλαμήδει ἐποίησε τὸν Οἴακα τὸν ἀδελφὸν Παλαμήδους ἐπιγράψαι εἰς τὰς ναῦς τὸν θάνατον αὐτοῦ, ἵνα φερόμεναι αὗται ἔλθωσιν εἰς τὸν Ναύπλιον τὸν πατέρα αὐτοῦ καὶ ἀπαγγείλωσι τὸν θάνατον αὐτοῦ.\\
\francais{Euripide, dans le \italique{Palamède}, a représenté Oiax, frère de Palamède, gravant sur des navires [le récit de] sa mort, afin qu'ils soient portés jusque chez Nauplios et lui apprennent sa mort.}}}
\comm{770}{ὡς ἐκεῖνος,}{Œax, frère de Palamède.}
\comm{772}{πόθεν οὖν γένοιντ᾽ ἄν μοι πλάται. πόθεν;}{Le vers est incomplet. Scaliger proposa la solution la plus instinctive : redoubler le πόθεν en fin de vers : πόθεν ; <πόθεν;> Cette solution est vraisemblablement la bonne. \ \ \ 1° \ Une simple haplographie l'explique. \ \ \ 2° \ Elle est une nouvelle occurrence des railleries artiphanesques quant à la manie du redoublement chez les tragiques. Le double πόθεν se trouve d'ailleurs dans l'\italique{Ion} (vers 1453 sq.) :\citComm{ἰὼ γύναι, πόθεν πόθεν ἔλαβες ἐμὸν\\
βρέφος ἐς ἀγκάλας;}Même si ce qui faisait vraiment battre le cœur d'Euripide, c'était le redoublement du ποῦ (\italique{Électre}, v. 487 sq.) :\citComm{ποῦ ποῦ νεᾶνις πότνι᾽ ἐμὴ δέσποινά τε,\\
Ἀγαμέμνονος παῖς, ὅν ποτ᾽ ἐξέθρεψ᾽ ἐγώ;}Dans les \italique{Oiseaux} (vers 921), c'est le redoublement de πάλαι qui est moqué : l'ennuyeux poète qui débite\citComm{πάλαι, πάλαι δὴ τήνδ᾽ ἐγὼ κλῄζω πόλιν}ne pénétrera pas plus dans Νεφελοκοκκυγίᾳ qu'il n'aura de chance de demeurer dans la Cité de Platon.}
\comm{773–4}{τί δ᾽ ἂν εἶτα διὰ τἀγάλματ᾽ ἀντὶ τῶν πλατῶν / γράφων διαρρίπτοιμι;}{Willems (\nieme{22} note, reproduite en annexe) a expliqué en quel sens il fallait prendre \grec{τὰ ἀγάλματα} : \grec{πᾶν ἐφ᾽ ᾧ τις ἀγάλλεται} (scholie à notre vers), \italique{tout ce qu'on offre et consacre aux dieux}. En l'occurrence, ce sont les \grec{πόπανα} que Mnésiloque a apportés (au vers 285) et non pas des \italique{ex-voto} comme on en voit encore apparaître dans les modernes traductions\nc{Après avoir encombré la scène d'une statue d'Apollon permanente d'après le vers 748, les commentateurs rajoutent ici des tablettes votives pour le décor.}.\\
Y a-t-il ici une allusion au proverbe grec ἄλλο γένος κώπης, «~un autre genre de rame~»\nc{Rapporté par Plutarque, \italique{Proverbes alexandrins}, 2, 52. Ou aussi dans \italique{Zu den Paroemiographen} de Cohn, p.42. Érasme l'a incorporé à ses \italique{Adages} (n°1621).} ?}
\comm{781}{τουτὶ τὸ ῥῶ μοχθηρόν·}{« Voilà un `Rhô' défectueux ! »... telle est la leçon de notre unique manuscrit et telles aussi les traductions embarrassées. Car la question fuse : pourquoi Mnésiloque nous parle-t-il d'un `ρ' ? Le scholiaste toujours prompt à lancer une ânerie propose d'y voir la troisième lettre du destinataire, Εὐρῑπίδη... ce qui ne répond pas à la question : pourquoi la troisième lettre ? Passons sur les autres explications du même acabit (si Aristophane avait un peu plus de glosateurs, tous les mots du dictionnaire comportant un rhô y seraient effectivement passés : Εὐριπίδη, χώρει, χαίρειν, \&c.)
Lisez simplement : \citComm{τουτὶ τορῶ μοχθηρόν.}L'aoriste τορεῖν est homérique et fait partie d'un groupe dialectal, également homérique, d'aoristes radicaux thématiques où s'est dégagé un vocalisme en `o' : ἔμολον de βλώσκω, ἔθορον de θρώσκω et ἔτορον de τείρω. Ces homérismes se sont introduits chez les Tragiques, et peuvent aussi se rencontrer dans certains passages d'Aristophane. On l'avait déjà noté pour ἔμολον : on peut maintenant y ajouter ἔτορον\nc{Je dois cette remarque à Jean-Victor Vernhes qui, entre autres qualités, est un très fin connaisseur de la langue homérique. Que le lecteur n'hésite pas à consulter le Dictionnaire étymologique de la langue grecque de Pierre Chantraine à l'entrée τορεῖν pour plus de détails.}. Hésychius (τ 1162) -- qui nous fut d'un grand secours comme à son habitude – l'avait déjà relevé :\citComm{τορεῖν· τορῆσαι, τρῆσαι, τεμεῖν ;}le verbe signifie \italique{percer}, \italique{perforer}, \italique{couper}.
Reconstituons la scène : Mnésiloque cherche un expédient pour envoyer son message. Heureusement, il connaît son euripide par cœur : il décide de procéder à la façon d'Œax dans le \italique{Palamède}. Manque de chance : il manque les rames (ἀλ ̓ οὐ πάρεισιν αἱ πλάται) pour graver dessus son message !.. qu'à cela ne tienne : les ustensiles de théâtres pour figurer les πόπανα feront bien l'affaire (ils sont de bois, après tout !) ; voilà donc nos morceaux de bois promus en tablettes (δέλτοι). Mnésiloque peut maintenant se mettre à graver la “cire” avec son stylet – un stylet ? quel stylet ? Nouveau détournement : il appelle `stylet' le coutelas avec lequel il vient d'égorger l'outre de vin. Seulement, à force de prendre des vessies pour des lanternes, Mnésiloque se fait rattraper par la dure sylvestre réalité : graver un message au couteau sur du bois se révèle bien plus difficile qu'avec un stylet sur une tablette de cire ; cela explique son cri de découragement : τουτὶ τορῶ μοχθηρόν, « ceci que je perce [il ne précise pas quoi, et pour cause : c'est à la fois un πόπανον, une πλάτη et un δέλτος]... est de mauvaise qualité ! »}
\comm{782}{χώρει, χώρει.}{Mnésiloque est en train de graver son message : χώρει, «~accours !~», enjoint-il à Euripide ; χώ...ρει χώ...ρει, «~Ac...cours ... Ac...cours~», se dicte-t-il à lui-même. Il l'écrit sur deux rames différentes, comme le pluriel αἱ πλάται le laissait présager et comme c'était aussi le cas dans la tragédie originale ; le scholiaste de notre \italique{Ravennas} précise bien : \citComm{ὁ γὰρ Οἴαξ ἐγχαράττει \textbf{πολλαῖς πλάταις} τὰ περὶ τὸν Παλαμήδην καὶ ἀφίησιν εἰς θάλασσαν, ὥστε μιᾷ γέ τινι τὸν Ναύπλιον περιπεσεῖν.\\
Œax grave effectivement de \italique{nombreuses rames} avec les circonstances touchant à $[$la mort de$]$ Palamède et les lance vers la mer de sorte que Nauplios tombe sur l'une d'elles.}À propos des ἀγάλματα, un scholiaste avait aussi écrit :\citComm{πίνακες γὰρ ἦσαν ἐν τῷ ἱερῷ, καὶ ἀπὸ τούτων λαβὼν γράφει, καὶ ῥίπτει αὐτὰς λέγων· `ἀπέλθετε, σημάνατε Ἐὐριπίδῄ.\\
C'était des tablettes dans le temple. Les prenant, il écrivit puis les lança, en disant : «~Allez ! Informez Euripide !~»}}
\comm{782–4}{ποίαν αὔλακα / βάσκετ᾽ ἐπείγετε πάσας καθόδους / κείνᾳ ταύτᾳ· ταχέως χρή.}{Les différents traducteurs voient dans la formule `ποίαν αὔλακα' une exclamation d'auto-contentement de Mnésiloque : «~Quelle belle lettre ! Quel sillon !~» Αὖλαξ est en effet attesté dans l'\italique{Anthologie palatine} (6, 68) au sens de sillon creusé par le stylet, c'est-à-dire le \italique{trait d'écriture}. Soit. Mais cette interprétation confère à ποίαν un sens qu'il n'a pas. Cet adjectif interrogatif sert à demander des précisions : \italique{de quelle nature ?} \italique{de quel genre ?} (Aristophane l'emploie fréquemment dans ses dialogues pour débrouiller -- ou souligner ? -- un quiproquo.) Il est utilisé en reprise d'un terme, reprise directe : Ἀγάθων ὁ κλεινὸς --- Ποῖος οὗτος Ἁγάθων; (vers 30) ; Τὸν δεῖνα γιγνώσκεις --- Τὸν δεῖνα; ποῖον; (vers 621) ;  Πρωτέως τάδ᾽ ἐστὶ μέλαθρα --- Ποίου Πρωτέως; (vers 874) ou implicite : Πρὸς ποῖον ἔργον; (vers 586) ; Χρείᾳ δὲ ποίᾳ τόνδ᾽ ἐπεισφέρεις λόγον; (vers 1163). Mnésiloque mentionne pour la première fois l'αὖλαξ ; mais reprise il y a bien\nc{Et si l'on en doutait, l'on pourrait toujours corriger le texte -- ce dont on se gardera bien ! Lire ποῖ ποτ' ἂν αὔλακα revient au même quant à l'action et dispense de supposer que l'αὖλαξ était le terme employé par Euripide. (ποτε \italique{gratia metri} mais dont la disparition relèverait d'une haplographie bien explicable. Au reste, la collocation `ποῖ ποτε' est très fréquente.)} : le mot devait sans aucun doute apparaître dans la tragédie d'Euripide. 
Où en est Mnésiloque de l'action ? Il vient de graver en deux exemplaires l'appel au secours, χώρει, qu'il compte faire parvenir à Euripide... comment ? Assurément, cette αὖλαξ désigne l'expédient qu'Œnax avait trouvé pour emporter ses messages.
On peut raisonnablement penser qu'il les déposait dans un \italique{cours d'eau} (qui envoyait ensuite les πλάται où l'une d'entre elles finissait par arriver à bon port, comme l'indique le scholiaste. Ça ne paraît pas très vraisemblable, mais la vraisemblance et une pièce d'Euripide font deux.)
C'est bien en ce sens qu'il faut prendre αὖλαξ, le \italique{sillon}, et, par extension, la \italique{rigole} ou le \italique{ruisseau}\nc{Voir le \LSJ \ et le \espagnol{Diccionario Griego-Español} IV à l'entrée αὖλαξ.}.\\
Mnésiloque ne s'embarrasse pas vraiment du lièvre qu'il vient de soulever : la question n'est là que pour souligner les invraisemblances de la mise en scène, comme Aristophane aime bien les relever. Il décide d'aller au plus rapide : il jette avec force les deux πλάται. (Une de chaque côté, vers chaque εἴσοδος, comme le laisse soupçonner πάσας καθ᾽ ὁδοὺς.) Cette \italique{nouveauté} (κείνᾳ ταύτᾳ) par rapport au modèle original d'Euripide a l'avantage d'accélérer l'action ; c'est ainsi qu'il faut comprendre le vers 784 : κείνᾳ ταύτᾳ· ταχέως χρή.}
\comm{785–845}{Parabase.}{La parabase des \italique{Thesmophories}, à l'instar de \italique{Lysistrata} et des \italique{Grenouilles}, ne présente ni ᾠδή/ἀντῳδή ni ἀντεπίρρημα\nc{Mon but n'est pas de proposer une description exhaustive de la structure des comédies d'Aristophane, ni même des \italique{Thesmophories} en particulier. Ce travail a déjà été très bien fait par White, \italique{The Verse of Greek Comedy} (1912). On lira aussi avec un grand intérêt les travaux de Jean Irigoin regroupés dans son anthologie posthume \italique{Le poète grec au travail} (2009) et notamment son article «~Prologue et \italique{parodos} dans la comédie d'Aristophane (\italique{Acharniens}, \italique{Paix}, \italique{Thesmophories})~» :  «~[Aristophane], ce poète comique à l'imagination débridée, s'astreignait, dans certaines parties de ses pièces, à respecter des règles concernant le choix des vers, leur disposition, et même leur nombre~». Le lecteur intéressé se reportera à ces travaux -- et à nombre d'autres -- auxquels je ne saurais me substituer.}.}
\comm{788}{ἔριδες νείκη στάσις ἀργαλέα}{Le vocabulaire est Homérique( \italique{Iliade}, XVII, v. 384-5) :\citComm{τοῖς δὲ πανημερίοις ἔριδος μέγα νεῖκος ὀρώρει / ἀργαλέης·\\
\francais{Durant tout le jour se prolongea cette lutte sanglante et terrible.}}}
\comm{789}{εἰ καὶ κακόν ἐσμεν, τί γαμεῖθ᾽ ἡμᾶς, εἴπερ ἀληθῶς κακόν ἐσμεν}{La perfidie d'Aristophane est à son comble quand il donne la parole aux femmes dans la parabase où elles assurent leur propre défense. « Si nous sommes un fléau, disent-elles aux hommes, pourquoi nous épousez-vous ? » Mais le spectateur, qui connaît, depuis l'école, son Hésiode par cœur et qui a déjà vu bien des comédies, répond intérieurement, comme mu par un ressort : « Parce qu'on n'échappe pas si facilement au piège tendu par Zeus, parce que la femme est un fléau nécessaire pour la perpétuation de l'espèce\nc{Hésiode, les \italique{Travaux et les jours}, v. 83-84 ; Susarion, fr. 1, 3-4 Kassel-Austin. Voir aussi ce que dit Renzo Tosi (§284 et §285 de son monumental \italique{Dictionnaire des sentences latines et grecques}, 2010) : Κακὸν ἀναγκαῖον γυνή ou \latin{Femina dulce malum} : la femme est un καλὸν κακόν (\italique{Théogonie}, v. 587). Quant au vers précédent, il y une allusion à peine voilée aux pandoriennes calamités déclenchées par la mythique -- et belle, surtout, pour nos malheurs ! -- Hélène. (Voir à ce propos Tosi, §350 : \latin{Nam fuit ante Helenam cunnus taeterrima belli causa}.) Il est donc tout naturel que Mnésiloqie réponde à cette parabase par une reprise de l'\italique{Hélène}.} ». Le poète comique ne fait ici, sous des apparences trompeuses, qu'entretenir le vieux réflexe de misogynie. Autre argument employé par les femmes pour montrer leur valeur : celui de la jalousie qu'elles excitent\nc{Aux vers 792-794.}. Tout commentaire est ici inutile... (Note d'Alain Blanchard\nc{Tirée de son remarquable article «~L'art de la mauvaise foi dans les \italique{Thesmophories} d'Aristophane~», in \italique{Dieux, héros et médecins grecs. Hommage à Fernand Robert.} (2001), p. 97-106. Tout amateur d'Aristophane ne peut pas ne pas aller prestement lire ces pages si riches (j'émets seulement quelques réserves sur le tout dernier paragraphe).
Blanchard y montre que « la mauvaise foi d'Aristophane s'est exercée, en ce qui concerne Euripide, sur un thème moins traditionnel que celui de la misogynie et qu'elle permet de mieux comprendre le sujet même de sa comédie des \italique{Thesmophories}. » De même que la défaite athénienne à Délion avait profité à la popularité de Socrate, qui s'était conduit héroïquement dans la débâcle comme on le sait par Platon, et que c'est ce bonheur d'une réputation rehaussée dans le malheur athénien qui avait provoqué Aristophane à l'écriture vengeresse des \italique{Nuées} ; de même, les \italique{Thesmophories} furent écrites peu après la désastreuse expédition de Sicile : alors qu'Athènes vivait ses heures les plus sombres, Euripide en profita pour redorer son blason : on fit appel à lui pour écrire un chant funèbre en l'honneur des Athéniens morts à Syracuse (Plutarque nous en a conservé deux vers). 
Anecdote plus significative encore : certains Athéniens, réduits en esclavage, ne durent leur salut qu'à Euripide, en récitant de ses prologues, si prisés des Syracusains. Que la poésie d'Euripide pût sauver qui que ce soit ?.. alors qu'elle est justement l'une des causes de la dégénérescence athénienne !? voilà une nouvelle fumisterie qu'Aristophane mettra à mal tout au long des \italique{Thesmophories} : les différents détournements des tragédies d'Euripide se heurteront toutes au mur de l'incompréhension des geôliers de Mnésiloque. Une fois bien cerné le contexte historique, l'on comprend mieux Aristophane : ce « jeu de massacre » n'est motivé que par la passion -- brûlante chez Aristophane -- de la justice : il cogne, cogne, et cogne encore contre tous ces charognards des misères athéniennes. \anglais{Now is a time to storm !} Pour conclure : si Euripide est misogyne, c'est là son moindre défaut.}.)}
\comm{797}{τὸ κακὸν ζητεῖτε θεᾶσθαι}{Suivant la belle correction de Kaye : τὸ κακὸν ζητεῖτε θεᾶσθαι, v. 797, au lieu de : ζητεῖ τὸ κακὸν τε θεᾶσθαι. (Note de Willems.)}
